危机情境是职业路径中最能体现形象韧性的环节。与“曝光下降”相比，更值得关注的是叙事权是否能够回到个体手中。危机后的恢复并不依赖单一声明，而需要在时间、平台与内容三方面同步重建。

\subsection{沉默、修复与重启}

本研究将危机后的路径拆解为三个连续阶段：沉默期负责降低扩散速度，修复期负责恢复可信度，重启期则负责形成新的内容节奏。只有当三者衔接得当，危机才可能转化为新的职业起点。

\subsection{比较性解释}

危机并不会自动导致职业衰退，它真正改变的是“谁来解释事件、通过什么平台解释、解释节奏由谁控制”。这也是为什么自有内容渠道在平台化时代变得重要：它不仅提供曝光，更提供解释权。

\subsection{阶段结果}

当职业主体能够在重启期持续输出时，公众对其职业身份的理解会从“被动回应争议”转回“主动定义新阶段”。这一机制与网络化公共空间中的持续自我呈现逻辑相符\cite{boyd2010,Goffman1959}。
